% GENERATED FILE. Edit canonical lessons, then run npm run generate:books.
% canonical-source-hash: fnv1a64:30fd3cf92737b48d
% canonical-lessons: PA-C151-kahina, PA-C151-jiuna, PA-C151-kharcna, PA-C151-bacauna, PA-C151-bijna

\chapter{To say, To live, To spend, To save, To sow}
\label{ch:pa-to-say-to-live-to-spend-to-save-to-sow}
\hlchaptermodality{151}

\begin{chapteropening}
\textbf{By the end of this chapter:} \emph{I can say to say, to live, to spend, to save and to sow.}

Complete the last lesson of chapter 151: I can say to say, to live, to spend, to save and to sow.
\end{chapteropening}

\section[\texorpdfstring{\pa{ਕਹਿਣਾ} (kahiṇā)}{ਕਹਿਣਾ (kahiṇā)}]{\pa{ਕਹਿਣਾ} (kahiṇā) — to say}

\label{lesson:PA-C151-kahina}



\begin{quote}
Before the new one: say the Punjabi for to wake up, then the Punjabi for to drive.
\end{quote}

\subsection*{You'll want to know: \pa{ਕਹਿਣਾ}}
\textbf{\pa{ਕਹਿਣਾ}} — \emph{kahiṇā} — \textquotedblleft{}to say\textquotedblright{}.

\begin{grammarlens}[title={What you've built}]
One more word for this chapter.
\end{grammarlens}

\subsection*{Your turn}
\begin{itemize}
\raggedright
  \item \emph{Say it:} \emph{kahiṇā}
  \item \emph{Say it:} \emph{kahiṇā}, once more
  \item \emph{Say it:} say \emph{kahiṇā}
  \item \emph{Recall:} say the Punjabi for to wake up, then the Punjabi for to drive, then say \emph{kahiṇā} again
\end{itemize}

\begin{tcolorbox}[breakable,colback=teal!4,colframe=teal!35!black,title={Before you move on}]
What does \pa{ਕਹਿਣਾ} mean? (\textquotedblleft{}To say\textquotedblright{}.) Say it once more.
\end{tcolorbox}
\section[\texorpdfstring{\pa{ਜੀਉਣਾ} (jīuṇā)}{ਜੀਉਣਾ (jīuṇā)}]{\pa{ਜੀਉਣਾ} (jīuṇā) — to live}

\label{lesson:PA-C151-jiuna}



\begin{quote}
Before the new one: say the Punjabi for to drive, then the Punjabi for to say.
\end{quote}

\subsection*{You'll want to know: \pa{ਜੀਉਣਾ}}
\textbf{\pa{ਜੀਉਣਾ}} — \emph{jīuṇā} — \textquotedblleft{}to live\textquotedblright{}.

\begin{grammarlens}[title={What you've built}]
One more word for this chapter.
\end{grammarlens}

\subsection*{Your turn}
\begin{itemize}
\raggedright
  \item \emph{Say it:} \emph{jīuṇā}
  \item \emph{Say it:} \emph{jīuṇā}, once more
  \item \emph{Say it:} say \emph{jīuṇā}
  \item \emph{Recall:} say the Punjabi for to drive, then the Punjabi for to say, then say \emph{jīuṇā} again
\end{itemize}

\begin{tcolorbox}[breakable,colback=teal!4,colframe=teal!35!black,title={Before you move on}]
What does \pa{ਜੀਉਣਾ} mean? (\textquotedblleft{}To live\textquotedblright{}.) Say it once more.
\end{tcolorbox}
\section[\texorpdfstring{\pa{ਖ਼ਰਚਣਾ} (kharcṇā)}{ਖ਼ਰਚਣਾ (kharcṇā)}]{\pa{ਖ਼ਰਚਣਾ} (kharcṇā) — to spend}

\label{lesson:PA-C151-kharcna}



\begin{quote}
Before the new one: say the Punjabi for to say, then the Punjabi for to live.
\end{quote}

\subsection*{You'll want to know: \pa{ਖ਼ਰਚਣਾ}}
\textbf{\pa{ਖ਼ਰਚਣਾ}} — \emph{kharcṇā} — \textquotedblleft{}to spend\textquotedblright{}.

\begin{grammarlens}[title={What you've built}]
One more word for this chapter.
\end{grammarlens}

\subsection*{Your turn}
\begin{itemize}
\raggedright
  \item \emph{Say it:} \emph{kharcṇā}
  \item \emph{Say it:} \emph{kharcṇā}, once more
  \item \emph{Say it:} say \emph{kharcṇā}
  \item \emph{Recall:} say the Punjabi for to say, then the Punjabi for to live, then say \emph{kharcṇā} again
\end{itemize}

\begin{tcolorbox}[breakable,colback=teal!4,colframe=teal!35!black,title={Before you move on}]
What does \pa{ਖ਼ਰਚਣਾ} mean? (\textquotedblleft{}To spend\textquotedblright{}.) Say it once more.
\end{tcolorbox}
\section[\texorpdfstring{\pa{ਬਚਾਉਣਾ} (bacāuṇā)}{ਬਚਾਉਣਾ (bacāuṇā)}]{\pa{ਬਚਾਉਣਾ} (bacāuṇā) — to save}

\label{lesson:PA-C151-bacauna}



\begin{quote}
Before the new one: say the Punjabi for to live, then the Punjabi for to spend.
\end{quote}

\subsection*{You'll want to know: \pa{ਬਚਾਉਣਾ}}
\textbf{\pa{ਬਚਾਉਣਾ}} — \emph{bacāuṇā} — \textquotedblleft{}to save\textquotedblright{}.

\begin{grammarlens}[title={What you've built}]
One more word for this chapter.
\end{grammarlens}

\subsection*{Your turn}
\begin{itemize}
\raggedright
  \item \emph{Say it:} \emph{bacāuṇā}
  \item \emph{Say it:} \emph{bacāuṇā}, once more
  \item \emph{Say it:} say \emph{bacāuṇā}
  \item \emph{Recall:} say the Punjabi for to live, then the Punjabi for to spend, then say \emph{bacāuṇā} again
\end{itemize}

\begin{tcolorbox}[breakable,colback=teal!4,colframe=teal!35!black,title={Before you move on}]
What does \pa{ਬਚਾਉਣਾ} mean? (\textquotedblleft{}To save\textquotedblright{}.) Say it once more.
\end{tcolorbox}
\section[\texorpdfstring{\pa{ਬੀਜਣਾ} (bījṇā)}{ਬੀਜਣਾ (bījṇā)}]{\pa{ਬੀਜਣਾ} (bījṇā) — to sow}

\label{lesson:PA-C151-bijna}



\begin{quote}
Before the new one: say the Punjabi for to spend, then the Punjabi for to save.
\end{quote}

\subsection*{You'll want to know: \pa{ਬੀਜਣਾ}}
\textbf{\pa{ਬੀਜਣਾ}} — \emph{bījṇā} — \textquotedblleft{}to sow\textquotedblright{}.

\begin{grammarlens}[title={What you've built}]
One more word for this chapter.
\end{grammarlens}

\subsection*{Your turn}
\begin{itemize}
\raggedright
  \item \emph{Say it:} \emph{bījṇā}
  \item \emph{Say it:} \emph{bījṇā}, once more
  \item \emph{Say it:} say \emph{bījṇā}
  \item \emph{Recall:} say the Punjabi for to spend, then the Punjabi for to save, then say \emph{bījṇā} again
\end{itemize}

\begin{tcolorbox}[breakable,colback=teal!4,colframe=teal!35!black,title={Before you move on}]
What does \pa{ਬੀਜਣਾ} mean? (\textquotedblleft{}To sow\textquotedblright{}.) Say it once more.
\end{tcolorbox}
